% Einheit 3 – Normalform und p-q-Formel – Nr. 32–44
\setcounter{aufgabe}{31}
\einheitenkopf[e3]{Einheit 3 von 4 · Normalform und p-q-Formel}
\zweigzeile{Hier lernst du, eine quadratische Gleichung in die Normalform zu bringen und mit der p-q-Formel zu lösen · neu in diesem Jahr (an der Oberschule je nach Buch Kl. 9 oder 10) · P10 oft · baut auf: Wurzelziehen (Einheit 1), Ausmultiplizieren und Ordnen (Einheit 2), binomische Formeln}

\begin{aufgabe}{Ich erkenne, welche Form eine Gleichung hat und welcher Lösungsweg passt.}
Rechne nicht. Kreuze den passenden Lösungsweg an.

\sachtabelle{lcccc}{Gleichung & Wurzel ziehen & Nullprodukt & rückwärts rechnen & p-q-Formel}{a) $3x^2 = 27$ & $\square$ & $\square$ & $\square$ & $\square$ \\ b) $(x - 4)\cdot(x + 1) = 0$ & $\square$ & $\square$ & $\square$ & $\square$ \\ c) $(x + 2)^2 = 11$ & $\square$ & $\square$ & $\square$ & $\square$ \\ d) $x^2 + 5x - 14 = 0$ & $\square$ & $\square$ & $\square$ & $\square$ \\ e) $x^2 - 7x + 6 = 0$ & $\square$ & $\square$ & $\square$ & $\square$}
\end{aufgabe}

\begin{aufgabe}{Ich kann p und q mit Vorzeichen ablesen.}
Die Normalform heißt $x^2 + p\cdot x + q = 0$. Rechne nichts aus. Schreibe $p$ und $q$ mit Vorzeichen auf.
\begin{geruest}
\gz{a}{$x^2 + 9x + 4 = 0$}{\feld{p}\feld{q}}
\gz{b}{$x^2 - 5x + 2 = 0$}{\feld{p}\feld{q}}
\gz{c}{$x^2 + 2x - 8 = 0$}{\feld{p}\feld{q}}
\gz{d}{$x^2 - x - 20 = 0$}{\feld{p}\feld{q}}
\gz{e}{$x^2 - 7x = 0$}{\feld{p}\feld{q}}
\end{geruest}
\end{aufgabe}

\begin{aufgabe}{Ich erkenne, ob die Vorzahl von $x^2$ eins ist.}
Rechne nicht. Musst du vor der Formel teilen? Kreuze an und gib an, durch welche Zahl.
\begin{teile}
\teil $x^2 + 4x - 5 = 0$ \hfill \kreuz{nein} \quad \kreuz{ja, teilen durch} \leerfeld
\teil $2x^2 - 6x + 4 = 0$ \hfill \kreuz{nein} \quad \kreuz{ja, teilen durch} \leerfeld
\teil $-x^2 + 3x + 10 = 0$ \hfill \kreuz{nein} \quad \kreuz{ja, teilen durch} \leerfeld
\teil $5x^2 + 10x - 15 = 0$ \hfill \kreuz{nein} \quad \kreuz{ja, teilen durch} \leerfeld
\teil $x^2 - 3x = 0$ \hfill \kreuz{nein} \quad \kreuz{ja, teilen durch} \leerfeld
\end{teile}
\end{aufgabe}

\verfahren{Lösen mit der p-q-Formel}

\begin{aufgabe}{Ich kann eine Gleichung in Normalform mit der p-q-Formel lösen.}
Löse mit der p-q-Formel.
\rechenplatz{4}
\begin{gleichungsraster}[4]
\gl{x^2 + 6x + 8 = 0} & \gl{x^2 + 10x + 21 = 0} \\
\gl{x^2 + 12x + 20 = 0} & \gl{x^2 + 6x + 5 = 0} \\
\gl{x^2 - 2x - 8 = 0} & \\
\end{gleichungsraster}
\end{aufgabe}

\begin{aufgabe}{Ich kann eine Gleichung erst ordnen oder normieren und dann mit der Formel lösen.}
Bringe die Gleichung in die Normalform und löse sie.
\rechenplatz{2}
\begin{gleichungsraster}[4]
\gl{x^2 + 5 = 6x} & \gl{x^2 = 2x + 3} \\
\gl{2x^2 + 4x - 16 = 0} & \gl{-x^2 + 6x - 8 = 0} \\
\gl{2x^2 + 2x - 12 = 0} & \\
\end{gleichungsraster}
\end{aufgabe}

\begin{aufgabe}{Ich kann mit der Formel erkennen, ob eine Gleichung zwei, eine oder keine Lösung hat.}
Löse mit der p-q-Formel.
\rechenplatz{2}
\begin{gleichungsraster}[3]
\gl{x^2 - 6x + 9 = 0} & \gl{x^2 - 2x + 5 = 0} \\
\anweisung{Berechne nur den Wert unter der Wurzel. Er heißt Diskriminante. Gib an, wie viele Lösungen die Gleichung hat.}
\gl[2]{x^2 - 3x + 2 = 0} & \gl[2]{x^2 + 5x + 7 = 0} \\
\gl[2]{x^2 - 7x + 12{,}25 = 0} & \\
\end{gleichungsraster}
\end{aufgabe}

\begin{aufgabe}{Ich kann Lösungen mit Wurzel und als Näherungswert angeben.}
Löse mit der p-q-Formel. Gib die Lösungen mit Wurzel und auf zwei Stellen nach dem Komma gerundet an.
\rechenplatz{2}
\begin{gleichungsraster}[4]
\gl{x^2 - 4x + 1 = 0} & \gl{x^2 + x - 1 = 0} \\
\gl[5]{3x^2 - 6x - 3 = 0} & \\
\end{gleichungsraster}
\end{aufgabe}

\begin{aufgabe}{Ich kann mit einer Lösung die Probe machen.}
Löse die Gleichung und mache mit einer Lösung die Probe.
\rechenplatz{2}
\begin{gleichungsraster}[5]
\gl{x^2 - 2x - 24 = 0} & \gl{3x^2 - 12x - 15 = 0} \\
\end{gleichungsraster}
\end{aufgabe}

\verfahren{Gleichungen mit Klammern}

\begin{aufgabe}{Ich kann eine Gleichung mit Klammern erst auflösen, dann ordnen und lösen.}
Löse die Gleichung.
\rechenplatz{3}
\begin{gleichungsraster}[5]
\gl{x\cdot(x + 3) = 18} & \gl{(x + 1)^2 = 3x + 7} \\
\gl{(x - 2)\cdot(x + 4) = 3x + 4} & \\
\end{gleichungsraster}
\end{aufgabe}

\verfahren{Den Lösungsweg wählen}

\begin{aufgabe}{Ich kann eine Gleichung auf dem passenden Weg lösen.}
Löse die Gleichung auf einem geschickten Weg.
\begin{gleichungsraster}[3]
\gl{4x^2 = 100} & \gl{x^2 + 7x = 0} \\
\gl{(x - 4)^2 = 1} & \gl[4]{x^2 - 3x - 4 = 0} \\
\end{gleichungsraster}
\begin{teile}
\teil Für die Parabel mit $f(x) = -x^2 + 4x + 9$ sollen alle Stellen $x$ berechnet werden, an denen $f(x) = -3$ gilt. Berechne sie. (P10 2023 FOR) \hfill \feld{x_1} \quad \feld{x_2}
\end{teile}
\end{aufgabe}

\begin{aufgabe}{Ich finde den Fehler bei der p-q-Formel.}
Finde den Fehler, benenne ihn und korrigiere die Rechnung.
\begin{teile}
\teil Emma löst $3x^2 + 6x - 9 = 0$:
\rechnung{p &= 6, \quad q = -9 \\ x_{1,2} &= -3 \pm \sqrt{9 + 9} \\ &= -3 \pm \sqrt{18}}
\teil Felix löst $x^2 - 8x + 12 = 0$:
\rechnung{x_{1,2} &= -4 \pm \sqrt{16 - 12} \\ &= -4 \pm 2 \\ x_1 &= -2, \quad x_2 = -6}
\teil Nora löst $x^2 - 10x + 16 = 0$:
\rechnung{x_{1,2} &= 5 \pm \sqrt{25 + 16} \\ &= 5 \pm \sqrt{41}}
\end{teile}
\end{aufgabe}

\begin{aufgabe}{Ich kann Regeln der p-q-Formel begründen.}
Begründe.
\begin{teile}
\teil Warum hat eine Gleichung keine Lösung, wenn unter der Wurzel eine negative Zahl steht?
\teil Warum muss man $2x^2 - 10x + 12 = 0$ erst durch 2 teilen, bevor man die p-q-Formel benutzt?
\end{teile}
\end{aufgabe}

\begin{aufgabe}{Ich kann die Schnittpunkte einer Geraden mit einer Parabel berechnen.}
Die Grafik zeigt die Parabel zu $f(x) = x^2 - 2$ und die Gerade zu $g(x) = -2x + 1$.

\begin{ksys}[xmin=-4,xmax=3,ymin=-3,ymax=8,ablesen]
\parabel{1}{0}{-2}{f}
\gerade{-2}{1}{g}
\end{ksys}

\begin{teile}
\teil Setze $f(x) = g(x)$ und berechne die Schnittpunkte. \hfill \feld{S_1} \quad \feld{S_2}
\teil Lies die Schnittpunkte am Graphen ab und vergleiche mit a).
\teil Gegeben sind die Parabel mit $f(x) = (x - 1)^2 - 4$ und die Gerade mit $g(x) = 3x + 3$. Berechne die Koordinaten der Schnittpunkte. (P10 2022 FOR) \hfill \feld{S_1} \quad \feld{S_2}
\end{teile}
\end{aufgabe}
